\chapter{Base Tecnológica}
\label{chap:basetecnoloxica}

\lettrine{X}{xxx},\dots

{ \color{red} Aspectos a considerar
    \begin{itemize}
      \item Párrafo introduciendo de forma general los siguientes subapartados. Las tecnologías suelen aparecer agrupadas. A continuación se indica una posible agrupación, bastante habitual en proyectos tipo de aplicación web.
      \item El objetivo de este apartado es enumerar y realizar una descripción super breve (1 o 2 párrafos) cada uno de los lenguajes, tecnologías, herrameintas, \dots utilizados durante la realización de este TFG. Muy importante, en toda la memoria pero en esta parte esencial, añadir referencia a bibliografía contrastada para cada ítem comentado.
      \item En la sección \ref{chap:desenrolo}, página \pageref{chap:desenrolo} situaremos cada una de las tecnologías en el punto de la arquitectura adecuado en nuestro proyecto y si se considera necesario, es en ese apartado en el que se realizará una justificación del uso de dichas tecnologías.
    \end{itemize}
}

\section{Lenguajes}
\label{sec:linguaxes}

\subsection{Java}
\subsection{HTML}
\subsection{CSS}

\section{Frameworks y librerías}
\label{sec:frameworks}

\subsection{Core}
\subsection{Web}
\subsection{Pruebas}

\section{Protocolos}
\label{sec:protocolos}

\section{Herramientas de desarrollo}
\label{sec:ferramentas}

\section{Servidores de aplicaciones}
\label{sec:servidores}

\section{Sistemas de gestión de base de datos}
\label{sec:bbdd}

\section{Otros?}
\label{sec:otros}